\documentclass{../../nndlcourse}

\begin{document}

\title{实验5 - 网络优化与正则化}
\author{数据科学与大数据技术专业}
\date{}
\maketitle


% ============================================================
\section{实验目标}
% ============================================================
\begin{itemize}
    \item 理解参数初始化对深层网络训练稳定性和收敛速度的影响.
    \item 掌握零初始化、大随机初始化和 He 初始化的实现方式, 并能解释它们在 ReLU 网络中的差异.
    \item 掌握 L2 正则化和 Dropout 的基本思想, 能够在损失函数、前向传播和反向传播中正确加入正则化项.
    \item 使用梯度检验验证反向传播代码, 理解数值梯度、解析梯度和相对误差之间的关系.
    \item 手动实现 Mini-batch 梯度下降、Momentum、Adam 和学习率衰减, 比较不同优化方法的训练现象.
    \item 能够根据损失曲线、准确率、决策边界和调试记录分析模型是否欠拟合、过拟合或优化困难.
\end{itemize}


% ============================================================
\section{实验环境}
% ============================================================
\begin{itemize}
    \item 操作系统: Windows / Linux / macOS
    \item Python 3.7 及以上
    \item 主要依赖库: \texttt{numpy}、\texttt{matplotlib}、\texttt{scipy}、\texttt{scikit-learn}、\texttt{h5py}
    \item 推荐开发环境: Jupyter Notebook 或 VS Code
    \item 实验文件:
    \begin{itemize}
        \item \texttt{W1A1/Initialization.ipynb}
        \item \texttt{W1A2/Regularization.ipynb}
        \item \texttt{W1A3/Gradient\_Checking.ipynb}
        \item \texttt{W2A1/Optimization\_methods.ipynb}
    \end{itemize}
\end{itemize}

如本地环境缺少依赖, 可在终端中执行:
\begin{verbatim}
pip install numpy matplotlib scipy scikit-learn h5py jupyter ipykernel
\end{verbatim}
若使用 Conda, 可执行:
\begin{verbatim}
conda install numpy matplotlib scipy scikit-learn h5py jupyter ipykernel
\end{verbatim}


% ============================================================
\section{背景知识}
% ============================================================

\subsection{训练深层网络为什么不只是写出反向传播}

前面实验已经实现了全连接网络、卷积网络和循环网络的核心计算. 但真实训练中, 即使前向传播和反向传播公式正确, 模型也可能因为初始化不合适、过拟合、梯度实现错误或优化器选择不当而训练失败. 本实验集中处理这些``训练工程''问题.

可以把一次可复现的训练实验拆成四个层次:
\begin{enumerate}
    \item \textbf{初始化}: 决定训练从哪里开始, 影响信号和梯度在网络中的尺度.
    \item \textbf{正则化}: 决定模型是否更偏向简单、稳定和可泛化的解.
    \item \textbf{梯度检验}: 验证自己写的反向传播是否真的等价于损失函数的导数.
    \item \textbf{优化算法}: 决定参数如何利用当前梯度和历史梯度向更低损失移动.
\end{enumerate}
本实验的重点不是背诵公式, 而是把这些技巧放回训练现象中理解: 为什么损失不降、为什么测试集变差、为什么某个优化器更快, 以及如何用实验记录说明自己的判断.

\subsection{参数初始化}

神经网络训练从一组初始参数开始. 若同一层所有隐藏单元的权重完全相同, 它们在前向传播中得到相同激活, 在反向传播中得到相同梯度, 后续更新也会保持相同, 网络无法学习多样化特征. 这就是权重零初始化的问题. 偏置初始化为零通常可以接受, 因为权重的随机性已经打破了对称性.

若权重初始化过大, 各层线性输出可能迅速放大, 使激活函数进入饱和区, 或让损失和梯度剧烈震荡. 若权重初始化过小, 信号在深层网络中逐层衰减, 梯度也可能变得很小. 因此, 初始化需要控制每层信号的方差.

对使用 ReLU 激活的网络, 常用 He 初始化:
\begin{align}
    W^{[l]} \sim \mathcal{N}\!\left(0,\frac{2}{n^{[l-1]}}\right),
    \quad
    W^{[l]}=\texttt{np.random.randn}(\cdots)
    \sqrt{\frac{2}{n^{[l-1]}}}.
\end{align}
其中 $n^{[l-1]}$ 是上一层神经元数量. 分子中的 2 与 ReLU 会将约一半负激活置零有关, 它用于补偿激活方差的下降.

\subsection{正则化与泛化}

当模型在训练集上表现很好, 但在测试集上明显变差时, 通常说明模型发生了过拟合. 正则化的目标不是让训练集损失最低, 而是让模型学到更平滑、更稳定、更能泛化的规律.

L2 正则化在原始交叉熵损失外加入权重平方和惩罚:
\begin{align}
    J_{\mathrm{reg}} =
    J_{\mathrm{data}}
    + \frac{\lambda}{2m}\sum_l \left\|W^{[l]}\right\|_F^2 .
\end{align}
对应到反向传播, 每一层权重梯度需要额外加上:
\begin{align}
    dW^{[l]} \leftarrow dW^{[l]} + \frac{\lambda}{m}W^{[l]}.
\end{align}
偏置 $b$ 通常不做 L2 正则化, 因为偏置不直接控制输入特征的放大方向, 对模型复杂度的贡献也远小于权重矩阵.

Dropout 在训练阶段随机关闭部分隐藏单元. 若 \texttt{keep\_prob} 表示保留概率, 常见的 inverted dropout 写法为:
\begin{align}
    D^{[l]} &= \mathbb{1}\left(\mathrm{rand}(A^{[l]}.shape)<\texttt{keep\_prob}\right),\\
    A^{[l]} &\leftarrow \frac{A^{[l]}\odot D^{[l]}}{\texttt{keep\_prob}}.
\end{align}
除以 \texttt{keep\_prob} 是为了保持激活值的期望尺度不变. 测试阶段不能使用随机 Dropout, 应使用完整网络做预测.

\subsection{梯度检验}

反向传播代码容易出现维度、转置、符号或求和轴错误. 梯度检验使用中心差分近似某个参数的梯度:
\begin{align}
    \frac{\partial J}{\partial \theta_i}
    \approx
    \frac{J(\theta_i+\varepsilon)-J(\theta_i-\varepsilon)}
    {2\varepsilon}.
\end{align}
将所有参数的数值梯度 \texttt{gradapprox} 与反向传播得到的梯度 \texttt{grad} 比较:
\begin{align}
    \mathrm{difference}
    =
    \frac{\left\|\texttt{gradapprox}-\texttt{grad}\right\|_2}
    {\left\|\texttt{gradapprox}\right\|_2+\left\|\texttt{grad}\right\|_2}.
\end{align}
若差异小于 $2\times 10^{-7}$, 通常可以认为反向传播实现基本正确. 梯度检验计算代价很高, 只适合调试阶段使用, 不应放入每次训练迭代.

\subsection{优化算法}

普通梯度下降直接沿负梯度方向更新:
\begin{align}
    W^{[l]} \leftarrow W^{[l]}-\alpha dW^{[l]},
    \quad
    b^{[l]} \leftarrow b^{[l]}-\alpha db^{[l]}.
\end{align}

Mini-batch 梯度下降每次只使用一小批样本估计梯度, 在计算效率和梯度稳定性之间折中. Momentum 使用历史梯度的指数加权平均减少震荡:
\begin{align}
    v_{dW}^{[l]} &= \beta v_{dW}^{[l]} + (1-\beta)dW^{[l]},\\
    W^{[l]} &\leftarrow W^{[l]}-\alpha v_{dW}^{[l]}.
\end{align}

Adam 同时维护一阶矩 $v$ 和二阶矩 $s$, 并进行偏差修正:
\begin{align}
    v_t &= \beta_1 v_{t-1}+(1-\beta_1)g_t,\\
    s_t &= \beta_2 s_{t-1}+(1-\beta_2)g_t^2,\\
    \hat{v}_t &= \frac{v_t}{1-\beta_1^t},
    \quad
    \hat{s}_t = \frac{s_t}{1-\beta_2^t},\\
    \theta_t &= \theta_{t-1}
    - \alpha\frac{\hat{v}_t}{\sqrt{\hat{s}_t}+\varepsilon}.
\end{align}
Momentum 与 Adam 都不会改变目标函数本身, 它们改变的是参数走向目标函数低谷的路径.


% ============================================================
\section{实验步骤}
% ============================================================

\subsection*{实验5.1: 参数初始化实验(W1A1)}

\textbf{实验目标: }在二维数据集上训练三层全连接网络, 比较零初始化、大随机初始化和 He 初始化对收敛速度、损失曲线和决策边界的影响.

网络结构为:
\begin{align}
    \mathrm{LINEAR}\to\mathrm{ReLU}\to
    \mathrm{LINEAR}\to\mathrm{ReLU}\to
    \mathrm{LINEAR}\to\mathrm{Sigmoid}.
\end{align}
Notebook 已提供训练函数 \texttt{model()}, 同学们需要补全三种初始化函数, 并用实验现象解释初始化为什么重要.

\subsubsection*{5.1.1 零初始化}

\textbf{练习 1: \texttt{initialize\_parameters\_zeros}}

实现要点:
\begin{itemize}
    \item 对每一层 $l$, 使用 \texttt{np.zeros} 初始化 $W^{[l]}$ 和 $b^{[l]}$.
    \item 参数字典键名为 \texttt{W1}、\texttt{b1}、\texttt{W2}、\texttt{b2} 等.
    \item 观察训练结果, 说明为什么隐藏层神经元会保持相同表示.
\end{itemize}
分析时不要只写``准确率低'', 应明确指出: 权重完全相同会导致同一层神经元收到相同梯度, 从而无法打破对称性.

\subsubsection*{5.1.2 大随机初始化}

\textbf{练习 2: \texttt{initialize\_parameters\_random}}

实现要点:
\begin{itemize}
    \item 权重使用 \texttt{np.random.randn(...)*10} 生成, 这是故意放大的初始化.
    \item 偏置仍初始化为零.
    \item 每层开始前设置 \texttt{np.random.seed(l)}, 以保证测试结果可复现.
    \item 观察损失曲线是否出现较大震荡、初始损失偏高或收敛缓慢.
\end{itemize}
该实验的目的不是推荐大随机初始化, 而是让同学们看到激活和梯度尺度失控时训练会变得不稳定.

\subsubsection*{5.1.3 He 初始化}

\textbf{练习 3: \texttt{initialize\_parameters\_he}}

实现要点:
\begin{itemize}
    \item 权重缩放因子为 \texttt{np.sqrt(2. / layers\_dims[l-1])}.
    \item 偏置初始化为零.
    \item 比较 He 初始化与前两种初始化在收敛速度、最终损失和决策边界上的差异.
\end{itemize}
报告中应说明 He 初始化为什么适合 ReLU 网络, 而不是只给出代码.

\subsection*{实验5.2: 正则化实验(W1A2)}

\textbf{实验目标: }在二维数据上训练三层网络, 比较无正则化、L2 正则化和 Dropout 三种设置对过拟合的影响.

网络结构为:
\begin{align}
    \mathrm{INPUT}(2)\to\mathrm{LINEAR}(20)\to\mathrm{ReLU}\to
    \mathrm{LINEAR}(3)\to\mathrm{ReLU}\to
    \mathrm{LINEAR}(1)\to\mathrm{Sigmoid}.
\end{align}
本实验应重点比较训练集准确率、测试集准确率和决策边界复杂度. 若训练集准确率很高但测试集较低, 应从过拟合角度解释.

\subsubsection*{5.2.1 L2 正则化损失}

\textbf{练习 1: \texttt{compute\_cost\_with\_regularization}}

在原始交叉熵代价基础上加入 L2 正则化项:
\begin{align}
    \frac{\lambda}{2m}
    \left(
    \left\|W^{[1]}\right\|_F^2+
    \left\|W^{[2]}\right\|_F^2+
    \left\|W^{[3]}\right\|_F^2
    \right).
\end{align}
实现时从 \texttt{parameters} 中取出 \texttt{W1}、\texttt{W2}、\texttt{W3}, 使用 \texttt{np.sum(np.square(W))} 计算平方和.

\subsubsection*{5.2.2 L2 正则化反向传播}

\textbf{练习 2: \texttt{backward\_propagation\_with\_regularization}}

在标准反向传播结果上加入 L2 梯度:
\begin{align}
    dW^{[l]} \leftarrow dW^{[l]}+\frac{\lambda}{m}W^{[l]}.
\end{align}
\texttt{db1}、\texttt{db2}、\texttt{db3} 保持原有计算方式. 完成后使用 \texttt{lambd=0.7} 训练模型, 对比测试集准确率和决策边界.

\subsubsection*{5.2.3 Dropout 前向传播}

\textbf{练习 3: \texttt{forward\_propagation\_with\_dropout}}

在第 1、2 层 ReLU 激活之后应用 Dropout:
\begin{enumerate}
    \item 生成与激活同形状的随机矩阵 \texttt{D1}、\texttt{D2}.
    \item 将随机矩阵转换为布尔掩码: \texttt{D = (D < keep\_prob)}.
    \item 用掩码关闭激活: \texttt{A = A * D}.
    \item 使用 \texttt{A = A / keep\_prob} 保持期望尺度.
    \item 将 \texttt{D1}、\texttt{D2} 存入 \texttt{cache}, 供反向传播使用.
\end{enumerate}
注意前向传播和反向传播必须使用同一组 Dropout 掩码; 不能在反向传播时重新随机生成.

\subsubsection*{5.2.4 Dropout 反向传播}

\textbf{练习 4: \texttt{backward\_propagation\_with\_dropout}}

反向传播时使用前向传播保存的同一组掩码:
\begin{align}
    dA^{[l]} \leftarrow
    \frac{dA^{[l]}\odot D^{[l]}}{\texttt{keep\_prob}}.
\end{align}
完成后使用 \texttt{keep\_prob=0.86} 训练模型, 分析 Dropout 是否缓解过拟合. 报告中应说明 Dropout 为什么通常会降低训练集表现, 但可能改善测试集表现.

\subsection*{实验5.3: 梯度检验(W1A3)}

\textbf{实验目标: }实现一维和多维梯度检验, 用数值梯度验证反向传播实现, 并修复 Notebook 中故意保留的反向传播错误.

\subsubsection*{5.3.1 一维梯度检验}

\textbf{练习 1--3: \texttt{forward\_propagation}、\texttt{backward\_propagation}、\texttt{gradient\_check}}

对于 $J(\theta)=\theta x$:
\begin{itemize}
    \item \texttt{forward\_propagation(x, theta)} 返回 $J$.
    \item \texttt{backward\_propagation(x, theta)} 返回 $d\theta=x$.
    \item \texttt{gradient\_check(x, theta, epsilon)} 计算中心差分近似梯度, 并与反向传播梯度比较.
\end{itemize}
预期 \texttt{difference} 应远小于 $2\times 10^{-7}$.

\subsubsection*{5.3.2 多维梯度检验}

\textbf{练习 4: \texttt{gradient\_check\_n}}

对三层网络中的所有参数进行梯度检验:
\begin{enumerate}
    \item 使用 \texttt{dictionary\_to\_vector} 将参数字典展平为向量.
    \item 对参数向量的每个分量分别构造 $\theta^+$ 和 $\theta^-$.
    \item 调用 \texttt{forward\_propagation\_n} 得到 $J^+$ 和 $J^-$.
    \item 计算每个分量的数值梯度.
    \item 使用 \texttt{gradients\_to\_vector} 将反向传播梯度展平, 再计算相对误差.
\end{enumerate}
Notebook 中的 \texttt{backward\_propagation\_n} 故意包含错误. 请使用梯度检验定位并修复 \texttt{dW2} 和 \texttt{db1} 附近的错误, 直到 \texttt{difference} 小于阈值. 报告中必须写明修复前后的 \texttt{difference}, 以及错误来自哪一个反向传播公式.

\subsection*{实验5.4: 优化算法(W2A1)}

\textbf{实验目标: }在二维数据集上比较不同优化方法, 并观察学习率衰减对训练稳定性的影响.

\subsubsection*{5.4.1 普通梯度下降}

\textbf{练习 1: \texttt{update\_parameters\_with\_gd}}

实现普通梯度下降参数更新. 对每一层 $l$ 同步更新:
\begin{align}
    W^{[l]} \leftarrow W^{[l]}-\alpha dW^{[l]},
    \quad
    b^{[l]} \leftarrow b^{[l]}-\alpha db^{[l]}.
\end{align}

\subsubsection*{5.4.2 Mini-batch 构造}

\textbf{练习 2: \texttt{random\_mini\_batches}}

实现小批量构造:
\begin{itemize}
    \item 用 \texttt{np.random.permutation(m)} 生成样本随机排列.
    \item 对 $X$ 和 $Y$ 使用相同索引同步打乱.
    \item 按 \texttt{mini\_batch\_size} 切分完整批次.
    \item 若最后剩余样本不足一个完整批次, 也要形成最后一个 mini-batch.
\end{itemize}
实现时要特别注意本课程数据矩阵通常每一列是一个样本, 因此打乱和切片主要发生在第 2 个维度.

\subsubsection*{5.4.3 Momentum}

\textbf{练习 3--4: \texttt{initialize\_velocity}、\texttt{update\_parameters\_with\_momentum}}

先用全零矩阵初始化速度字典 \texttt{v}, 键名为 \texttt{dW1}、\texttt{db1} 等. 随后实现 Momentum 更新:
\begin{align}
    v_{dW}^{[l]} &= \beta v_{dW}^{[l]}+(1-\beta)dW^{[l]},\\
    v_{db}^{[l]} &= \beta v_{db}^{[l]}+(1-\beta)db^{[l]},\\
    W^{[l]} &\leftarrow W^{[l]}-\alpha v_{dW}^{[l]},\\
    b^{[l]} &\leftarrow b^{[l]}-\alpha v_{db}^{[l]}.
\end{align}
分析时可把 Momentum 理解为对梯度方向做平滑, 它通常能减少在狭长谷底中的来回震荡.

\subsubsection*{5.4.4 Adam}

\textbf{练习 5--6: \texttt{initialize\_adam}、\texttt{update\_parameters\_with\_adam}}

初始化一阶矩字典 \texttt{v} 和二阶矩字典 \texttt{s}. 每次更新时依次完成:
\begin{enumerate}
    \item 更新一阶矩 \texttt{v}.
    \item 更新二阶矩 \texttt{s}.
    \item 对 \texttt{v} 和 \texttt{s} 做偏差修正.
    \item 用修正后的矩估计更新参数.
\end{enumerate}
函数应返回 \texttt{parameters}、\texttt{v}、\texttt{s}、\texttt{v\_corrected}、\texttt{s\_corrected}. 注意时间步 $t$ 从 1 开始参与偏差修正, 不能把 $t=0$ 代入分母.

\subsubsection*{5.4.5 学习率衰减}

\textbf{练习 7--8: \texttt{update\_lr}、\texttt{schedule\_lr\_decay}}

实现两种学习率调度:
\begin{align}
    \alpha &= \frac{\alpha_0}{1+\texttt{decayRate}\cdot\texttt{epochNum}},\\
    \alpha &= \frac{\alpha_0}
    {1+\texttt{decayRate}\cdot
    \left\lfloor \texttt{epochNum}/\texttt{timeInterval}\right\rfloor}.
\end{align}
对比固定学习率、连续衰减和分段衰减在损失曲线上的差异. 分析时应说明衰减学习率为什么可能让后期训练更稳定, 也要注意过早衰减可能导致收敛停滞.


% ============================================================
\section{核心函数总结}
% ============================================================

\begin{center}
\begin{tabular}{|p{1.4cm}|p{5.1cm}|p{5.8cm}|}
\hline
\textbf{NB} & \textbf{函数} & \textbf{重点} \\
\hline
W1A1 & 零初始化 & 观察权重零初始化导致的对称性问题. \\
\hline
W1A1 & 大随机初始化 & 观察过大随机权重带来的训练不稳定. \\
\hline
W1A1 & He 初始化 & 掌握适用于 ReLU 网络的方差缩放初始化. \\
\hline
W1A2 & L2 损失 & 在损失中加入权重平方惩罚. \\
\hline
W1A2 & L2 反向传播 & 在权重梯度中加入 $\frac{\lambda}{m}W$. \\
\hline
W1A2 & Dropout 前向传播 & 生成、应用并保存 Dropout 掩码. \\
\hline
W1A2 & Dropout 反向传播 & 使用同一组掩码传递梯度并保持尺度. \\
\hline
W1A3 & \texttt{gradient\_check\_n} & 展平参数并逐分量计算数值梯度. \\
\hline
W2A1 & \texttt{random\_mini\_batches} & 同步打乱样本并切分 mini-batch. \\
\hline
W2A1 & Adam 参数更新 & 实现一阶矩、二阶矩、偏差修正和参数更新. \\
\hline
\end{tabular}
\end{center}


% ============================================================
\section{实验作业}
% ============================================================
\begin{enumerate}
    \item 完成 \texttt{W1A1/Initialization.ipynb}, 比较零初始化、大随机初始化和 He 初始化的损失曲线、准确率和决策边界.
    \item 完成 \texttt{W1A2/Regularization.ipynb}, 比较无正则化、L2 正则化和 Dropout 的训练集表现、测试集表现和决策边界.
    \item 完成 \texttt{W1A3/Gradient\_Checking.ipynb}, 写明初始 \texttt{difference}、发现的问题、修复位置和修复后的 \texttt{difference}.
    \item 完成 \texttt{W2A1/Optimization\_methods.ipynb}, 比较 Mini-batch GD、Momentum、Adam 以及学习率衰减对训练过程的影响.
    \item 至少选择一个实验结果图或表格进行解释: 例如损失曲线、决策边界、训练/测试准确率对比表.
    \item \textbf{选做}: 自行调整 \texttt{lambd}、\texttt{keep\_prob}、\texttt{mini\_batch\_size}、学习率或优化器超参数, 分析它们对训练过程的影响.
    \item \textbf{选做}: 尝试用 AutoResearch 的思路优化本实验中的某一个任务.
\end{enumerate}


% ============================================================
\section{提交要求}
% ============================================================

\textbf{提交前请逐项检查材料是否齐全.} 本次实验提交内容包括 Notebook 源文件和实验过程记录. 本课程作业上传只接受文本文件, 不接受 PDF、Word 文档、图片、截图、压缩包等二进制文件. 采用这一规则, 一方面是为了引入 AI 辅助评阅, 另一方面也是希望大家养成用清晰文本表达实验过程和结论的习惯.

\begin{enumerate}
    \item \textbf{Notebook 文件}
    \begin{itemize}
        \item 提交补全后的四个 Notebook:
        \begin{itemize}
            \item \texttt{Initialization.ipynb}
            \item \texttt{Regularization.ipynb}
            \item \texttt{Gradient\_Checking.ipynb}
            \item \texttt{Optimization\_methods.ipynb}
        \end{itemize}
        \item Notebook 本身属于文本格式, 可以上传. 请保留必要的运行结果、损失曲线、准确率、决策边界和关键结论.
        \item 如果实验报告必须包含配图, 请将图片保留在 Notebook 输出中, 形成完整实验报告; 不要单独上传图片、PDF 或压缩包.
        \item 若完成选做内容, 请在 Notebook 或实验记录中保留配置、指标、结果和简要说明.
    \end{itemize}
    \item \textbf{实验过程与 AI 互动记录文本文件(.md / .txt)}
    \begin{itemize}
        \item 必须额外提交一份文本文件, 作为实验过程与 AI 互动记录, 推荐使用 \texttt{.md} 或 \texttt{.txt}.
        \item 记录内容至少包括: 遇到的主要问题、查阅资料或询问 AI 的问题、自己的判断、调试尝试、修改思路和实验结论.
        \item 不要把 AI 的回复原文放入记录中, 也不要复制粘贴大段代码; 记录应只包含你自己写的内容.
        \item 若使用 \LaTeX{}、Word 或其他文档工具撰写报告, 请使用 pandoc 等工具转换为纯文本或 Markdown 后提交.
        \item 建议文件命名为 \texttt{AI问答记录.md} 或 \texttt{实验过程记录.md}.
    \end{itemize}
\end{enumerate}


% ============================================================
\section{关于作业的重要说明}
% ============================================================

本实验包含大量看似简单但实际容易出错的实现细节, 例如矩阵形状、Dropout 掩码保存、Adam 偏差修正和学习率调度的轮次计算. 完成作业时, 不仅要让测试单元通过, 还要能解释为什么这样实现.

对于必做内容, 可以使用 AI 工具辅助理解概念、解释报错、查询 API 用法或检查矩阵维度, 但不得直接让 AI 代写全部必做代码后不加理解地提交. 若使用 AI 工具, 必须在文本记录中如实写明自己的提问、自己的判断和后续修改过程.

AI 互动记录应只保留学生自己写的内容. 不要粘贴 AI 的完整回复, 不要粘贴复制来的大段代码, 也不要把聊天记录原样导出后提交. 记录的重点是说明你如何提出问题、如何判断回答是否可靠, 以及如何把这些信息转化为自己的实现和分析.

选做内容可以更自由地使用 AI 工具, 但仍建议记录关键尝试过程和实验结论.

\end{document}
